\section{결론 및 시사점}

본 연구는 한국산업은행을 정책금융기관으로서의 관점에서 분석하고, 그 역할과 한계를 구조적으로 검토하였다.

한국산업은행은 설립 이후 산업화, 금융자유화, 외환위기, 그리고 최근의 산업 구조 전환에 이르기까지 각 시기마다 서로 다른 역할을 수행해왔다. 이러한 변화는 산업은행이 고정된 기능을 수행하는 기관이 아니라, 경제 환경과 정책 필요에 따라 기능이 조정되는 \textbf{적응적 제도(adaptive institution)}임을 보여준다.

특히 산업은행은 다음과 같은 세 가지 핵심 기능을 수행한다:

\begin{itemize}
    \item 시장실패를 보완하는 자금 공급 기능
    \item 경기 변동을 완화하는 거시경제 안정화 기능
    \item 산업 구조를 형성하는 정책금융 기능
\end{itemize}

이러한 기능은 민간 금융만으로는 달성하기 어려운 경제적 목표를 가능하게 하며, 산업은행이 국가 경제 시스템에서 중요한 위치를 차지하는 이유를 설명한다.

그러나 동시에 산업은행은 구조적 한계를 내포하고 있다.

첫째, 정부 지배구조로 인해 정치적 개입 가능성이 존재하며, 이는 자원 배분의 효율성을 저해할 수 있다.  
둘째, 반복적인 정책금융 지원은 기업의 도덕적 해이를 유발하고 시장 규율을 약화시킬 수 있다.  
셋째, 민간 금융기관과의 경쟁 관계는 금융시장 내 역할 중복과 왜곡 문제를 발생시킬 수 있다.

이러한 문제는 다음과 같은 근본적인 딜레마로 요약된다:

\begin{center}
\textbf{공공성 (policy objective) vs 효율성 (market efficiency)}
\end{center}

즉, 산업은행은 정책 목표를 달성하기 위해 시장에 개입해야 하지만, 그 개입 자체가 시장의 효율성을 저해할 수 있는 구조적 긴장을 가진다.

따라서 향후 정책금융의 방향은 이 두 요소 간의 균형을 어떻게 설정할 것인가에 달려 있다.

이를 위해 다음과 같은 정책적 시사점을 도출할 수 있다:

\subsection*{정책적 시사점}

\begin{itemize}
    \item \textbf{개입의 범위 명확화} \\
    시장이 실패하는 영역에 한정하여 정책금융을 집중하고, 민간 금융과의 역할 중복을 최소화할 필요가 있다.

    \item \textbf{거버넌스 개선} \\
    경영 독립성과 책임성을 강화하여 정치적 개입 가능성을 줄이고 의사결정의 효율성을 높여야 한다.

    \item \textbf{시장 규율 유지} \\
    구조조정 및 금융 지원 과정에서 시장 원칙을 유지하여 도덕적 해이를 방지해야 한다.

    \item \textbf{정책금융 체계 재정립} \\
    정책금융기관 간 기능 중복을 최소화하고, 효율적인 역할 분담 구조를 구축할 필요가 있다.
\end{itemize}

결론적으로, 한국산업은행은 단순한 금융기관이 아니라 국가 경제의 구조와 안정성을 동시에 조정하는 핵심 제도이다.

그러나 그 효과는 제도의 존재 자체가 아니라, \textbf{어떻게 설계되고 운영되는가}에 의해 결정된다.

따라서 향후 산업은행의 발전 방향은 정책금융의 필요성을 인정하면서도, 시장 메커니즘과의 조화를 이루는 방향으로 설정되어야 할 것이다.